\documentclass[10pt]{article}
\usepackage{aiguide}
\guidekicker{Guide 2 of 3}
\newcolumntype{L}{>{\raggedright\arraybackslash}X}
\newcommand{\toolrow}[4]{\textbf{#1}\newline{\footnotesize\href{#2}{#3}} & #4 \\ \addlinespace[0.15em]}
\newcommand{\tooltable}[1]{%
  \begin{tabularx}{\linewidth}{@{}p{0.31\linewidth}L@{}}
  #1
  \end{tabularx}}
\newcommand{\caution}[1]{{\small\color{muted}\textbf{Watch out:} #1}\par}
\begin{document}
\newgeometry{margin=2.1cm,top=1.9cm,bottom=2.1cm}
\makeguidetitle{Which AI Tool for Which Job}{A map of AI tools for students of mathematics, grouped by what you need to do.}

Pick the tool for the job, not the other way round. A chat assistant is good at
explaining and bad at arithmetic; a computer algebra system is the opposite.
Facts below come from each tool's official pages and were last checked in
November 2026. These tools change quickly, so check the page before you rely on
a feature.

\section{Explaining and tutoring: general assistants}
\tooltable{%
\toolrow{ChatGPT}{https://chatgpt.com/}{chatgpt.com}{General assistant from OpenAI. Study mode asks questions and gives hints instead of the final answer.}
\toolrow{Gemini}{https://gemini.google.com/}{gemini.google.com}{General assistant from Google. Guided Learning walks you through a problem step by step.}
\toolrow{Grok}{https://grok.com/}{grok.com}{General assistant from xAI (now SpaceXSI), with real-time answers from the web and X.}
\toolrow{Claude}{https://claude.com/}{claude.com}{General assistant from Anthropic.}
}
\caution{They write wrong answers as fluently as right ones. Age rules differ:
ChatGPT and Grok are 13+ with a parent's permission under 18, Gemini needs a Google Account
you may manage yourself (14+ in Italy), and Claude is 18+.}

\section{Studying from your own material}
\tooltable{%
\toolrow{Gemini Notebook\newline{\footnotesize\mdseries(formerly NotebookLM)}}{https://notebook.google/}{notebook.google}{Answers questions about the notes, slides, and PDFs you upload, cites the passages it used, and makes study guides, flashcards, and quizzes.}
}
\caution{Answers come only from what you upload, so a gap in your notes is a gap
in the answer. Open the cited passage before you trust a summary.}

\section{Computing and checking}
\tooltable{%
\toolrow{Wolfram|Alpha}{https://www.wolframalpha.com/}{wolframalpha.com}{Solves equations, limits, and integrals from plain-language input. Step-by-step solutions need Pro.}
\toolrow{GeoGebra CAS Calculator}{https://www.geogebra.org/cas}{geogebra.org/cas}{Free computer algebra in the browser: solve, factor, differentiate.}
\toolrow{Desmos Graphing Calculator}{https://www.desmos.com/calculator}{desmos.com/calculator}{Plot a function to check zeros, signs, and asymptotes.}
}
\caution{These compute exactly, but only what you typed. Check that the input is
the problem you meant.}

\section{University work: LaTeX, code, papers, proofs}
\tooltable{%
\toolrow{Overleaf AI features}{https://docs.overleaf.com/integrations-and-add-ons/ai-features}{docs.overleaf.com}{Error Assist explains LaTeX errors and suggests fixes. Free within daily limits.}
\toolrow{Gemini in Colab}{https://developers.google.com/colab}{developers.google.com/colab}{Python notebooks where Gemini generates, explains, and debugs code. Learn Mode guides you instead of writing the code.}
\toolrow{Semantic Scholar}{https://www.semanticscholar.org/}{semanticscholar.org}{Search research papers. TLDRs give a one-sentence AI summary of each result.}
\toolrow{Lean and Mathlib}{https://lean-lang.org/}{lean-lang.org}{Advanced: a proof assistant that checks every step of a proof, and its library of formalized mathematics (\href{https://mathlib.org/}{mathlib.org}).}
}
\caution{Read every line of generated LaTeX or code before you use it. A
one-sentence summary is not the paper.}

\section{Coding agents}
\tooltable{%
\toolrow{Grok Build}{https://x.ai/build}{x.ai/build}{xAI's coding agent for the terminal: plans, edits files, runs builds and tests. Free to try.}
\toolrow{Claude Code}{https://claude.com/product/claude-code}{claude.com/product/claude-code}{Anthropic's coding agent for the terminal, editors, and the web. Asks before it changes files or runs commands. Paid Claude plans.}
\toolrow{Codex}{https://chatgpt.com/codex/}{chatgpt.com/codex}{OpenAI's coding agent in ChatGPT, your editor, and the terminal, with cloud environments for each repository. Included in paid ChatGPT plans.}
\toolrow{Google Antigravity}{https://antigravity.google/}{antigravity.google}{Google's agentic development platform: an app, a CLI, an IDE, and editor extensions where agents write, run, and test code.}
}
\caution{They edit and run code on your machine or in your repository. Review
every change, use version control, and do not hand them graded coursework.}

\section{Agents and personal assistants}
Agents do not just answer: they carry out multi-step tasks such as browsing,
running code, organising files, and scheduling.

\tooltable{%
\toolrow{Grok Bot}{https://x.ai/bot}{x.ai/bot}{Agents with their own cloud computer that use apps and websites, write and run code, and work on a schedule.}
\toolrow{OpenAI Dots}{https://chatgpt.com/features/dots/}{chatgpt.com/features/dots}{Always-on ChatGPT agents with their own cloud computer that keep working on your projects between conversations. Paid plans, 18+, not yet available in the EU.}
\toolrow{Meta Muse}{https://about.fb.com/news/2026/09/introducing-muse-personal-ai-agent/}{about.fb.com}{Meta's personal agent on its own cloud computer: it browses, fills in forms, and books things, and asks before it sends an email or buys. Rolling out in the US, free with a usage limit.}
\toolrow{Claude Duet}{https://claude.com/product/claude-duet}{claude.com/product/claude-duet}{Anthropic's proactive agent on its own cloud computer that works in your connected apps and asks before it sends, pays, or deletes. US research preview for paid plans, by waitlist.}
}
\caution{Agents act on your accounts and can be confidently wrong. Review what
they do, keep personal data and school accounts out, and never let them submit
graded work.}

\section{Quick map}
\begin{tabularx}{\linewidth}{@{}p{0.42\linewidth}L@{}}
\toprule
\textbf{I want to\,\dots} & \textbf{Reach for} \\
\midrule
understand an idea or get a hint & a general assistant in study or tutor mode \\
revise from my own notes & Gemini Notebook \\
check an answer or a calculation & Wolfram|Alpha, GeoGebra CAS, Desmos \\
fix a LaTeX error & Overleaf Error Assist \\
understand or debug Python & Gemini in Colab \\
work on a coding project & a coding agent, inside version control \\
find papers on a topic & Semantic Scholar, then read the papers \\
be certain a proof is correct & Lean, or a careful human reader \\
hand off a multi-step chore & an agent, and check everything it did \\
\bottomrule
\end{tabularx}

\end{document}
